%%
%% ---------- header ----------------------------------------------------------
%%
%% project       kaneton
%%
%% license       kaneton
%%
%% file          /home/mycure/kane.../book/arch-ia32/virtual/ia32-virtual.tex
%%
%% created       julien quintard   [mon may 14 19:18:46 2007]
%% updated       julien quintard   [mon may 14 19:18:48 2007]
%%

%
% path
%

\def\path{../../..}

%
% template
%

%
% ---------- header -----------------------------------------------------------
%
% project       kaneton
%
% license       kaneton
%
% file          /home/mycure/kaneton/view/template/book.tex
%
% created       julien quintard   [wed may 16 18:16:50 2007]
% updated       julien quintard   [sun jun  3 12:00:37 2007]
%

%
% class
%

\documentclass[10pt,a4wide]{book}

%
% packages
%

\usepackage[english]{babel}
\usepackage[T1]{fontenc}
\usepackage{a4wide}
\usepackage{fancyheadings}
\usepackage{multicol}
\usepackage{indentfirst}
\usepackage{graphicx}
\usepackage{color}
\usepackage{xcolor}
\usepackage{verbatim}

\usepackage{aeguill}

\usepackage[Lenny]{\path/package/fncychap}

\pagestyle{fancy}

%
% -
%

\renewcommand{\-}{\vspace{\parskip}}

%
% toc
%

\newcommand{\toc}
  {
    \tableofcontents
    \setlength{\footrulewidth}{0.3pt}
    \setlength{\parindent}{0.3cm}
    \setlength{\parskip}{2ex plus 0.5ex minus 0.2ex}
  }

%
% logos
%

\newcommand{\logos}
  {
    \begin{center}
      \includegraphics[scale=0.8]{\path/logo/kaneton.pdf}
    \end{center}
  }

%
% colors
%

\definecolor{functioncolor}{rgb}{0.40,0.00,0.00}
\definecolor{commandcolor}{rgb}{0.00,0.00,0.40}
\definecolor{verbatimcolor}{rgb}{0.00,0.40,0.00}
\definecolor{noticecolor}{rgb}{0.87,0.84,0.02}

%
% function
%

\newcommand\function[3]{
  \begin{tabular}{p{0.2cm}p{13.8cm}}
  & {\color{functioncolor}\textbf{#1}}#2
  \end{tabular}

  \begin{tabular}{p{1cm}p{13cm}}
  & #3
  \end{tabular}}

%
% align
%

\newcommand\align[1]{
  \\ & \hspace{#1}}

%
% argument
%

\newcommand\argument[1]{\textit{#1}}

%
% command
%

\newcommand\command[2]{
  \begin{tabular}{p{0.2cm}p{13.8cm}}
  & {\color{commandcolor}\textbf{#1}}
  \end{tabular}

  \begin{tabular}{p{1cm}p{13cm}}
  & #2
  \end{tabular}}

%
% notice
%

\newcommand\notice[1]{
  {\color{noticecolor}\textbf{Notice}}

  \begin{tabular}{p{0.2cm}p{13.8cm}}
  & #1
  \end{tabular}}

%
% subsubsubsection
%

\newcommand\subsubsubsection[1]{\textbf{#1}}

%
% example
%

\newcommand\example[1]{
  \textit{Example:}

  \begin{tabular}{p{0.2cm}p{13.8cm}}
  & \textit{#1}
  \end{tabular}}

%
% warning XXX
%

\renewcommand{\familydefault}{\sfdefault}

%
% verbatim stuff
%

\makeatletter

\renewcommand{\verbatim@font}
  {\ttfamily\footnotesize\selectfont}

\def\verbatim@processline{
  \hskip15ex{\color{verbatimcolor}\the\verbatim@line}\par
}

\makeatother

%
% header
%

\rhead{}
\rfoot{\scriptsize{The kaneton microkernel project}}

\date{\scriptsize{\today}}


%
% header
%

\lhead{\scriptsize{The kaneton microkernel :: ia32/educational architecture}}
\rhead{}

%
% title
%

\title{The kaneton microkernel :: ia32/educational architecture
       \logos}

%
% authors
%

\author{\small{Matthieu Bucchianeri} and
        \small{Renaud Voltz}}

%
% document
%

\begin{document}

%
% title
%

\maketitle

%
% --------- text --------------------------------------------------------------
%

This document describes the kaneton microkernel reference project over
Intel IA-32 architecture.

This document should be used by every student willing implement the
kaneton microkernel as people looking for more details on the kaneton
microkernel implementation on Intel architecture.

All the kaneton documents are available on
the official website
  \footnote{http://www.kaneton.org}.

\textbf{Prerequisites}: this document details the kaneton implementation
for IA-32 architecture. A knowledge of this architecture's basics is
strongly recommended.

%
% toc
%

\tableofcontents

%
% chapters
%

%%
%% licence       kaneton licence
%%
%% project       kaneton
%%
%% file          /home/mycure/kaneton/view/papers/kaneton/overview.tex
%%
%% created       matthieu bucchianeri   [mon jan 30 17:09:45 2006]
%% updated       julien quintard   [thu mar  2 13:12:22 2006]
%%

%
% overview
%

\chapter{Overview}

XXX ce chapitre va vous aider a reconnaitre les fonctionnalites principale
XXX d'un kernel dans kaneton.

The kaneton microkernel is only the core of an operating system.
Main tasks like hardware drivers or user services are implemented as
\textbf{servers}. So the microkernel only has a few functionalities to
provide:

\begin{itemize}
  \item
    Memory management.
  \item
    Process management.
  \item
    Communication.
  \item
    Events.
\end{itemize}

In this chapter we will describe briefly these tasks and all the
associated managers.

%
% memory management
%

\section{Memory Management}

Handling the memory -- from virtual address space to physical
addressing -- is done by three major managers, the \textbf{as},
\textbf{segment} and \textbf{region} managers.

%
% as
%

\subsection{as}

The address space manager just manages the different address spaces
used by the kaneton tasks.

In kaneton, we call an \textbf{as - address space} a list of memory
locations referenced by a task. Each task has its own address space.

%
% segment
%

\subsection{segment}

The segment manager just manages the segments reserved by
the different kaneton entities including the kernel, the drivers etc..

In kaneton terms a \textbf{segment} is a contiguous area of reserved
physical memory.

%
% region
%

\subsection{region}

The region manager keeps track of regions used to map segments for
each address space reserved on the system.

In kaneton, a \textbf{region} is contiguous area of virtual memory
mapping a segment's part.

%
% process management
%

\section{Process Management}

XXX

%
% communication
%

\section{Communication}

XXX

%
% events
%

\section{Events}

XXX

%%
%% licence       kaneton licence
%%
%% project       kaneton
%%
%% file          /home/buckman/kaneton/view/books/ia32-virtual/bootloader.tex
%%
%% created       matthieu bucchianeri   [sat sep  2 11:40:18 2006]
%% updated       matthieu bucchianeri   [sat sep  2 11:40:32 2006]
%%

%
% bootloader
%

\chapter{Bootloader}

%
% ---------- header -----------------------------------------------------------
%
% project       kaneton
%
% license       kaneton
%
% file          /home/enguerrand/...view/book/architecture/mips64/memory.tex
%
% created       enguerrand raymond   [tue oct  7 08:54:35 2008]
% updated          [tue apr 28 14:30:30 2009]
%

%
% memory
%

\chapter{Memory management}

This chapter describes the memory management developped for MIPS64 microprocessors.

\subsection{Bootloader}

\subsection{kaneton kernel}

%%
%% licence       kaneton licence
%%
%% project       kaneton
%%
%% file          /home/buckman/kaneton/view/books/ia32-virtual/event.tex
%%
%% created       matthieu bucchianeri   [sat sep  2 11:40:54 2006]
%% updated       matthieu bucchianeri   [sun sep  3 12:24:43 2006]
%%

%
% event management
%

\chapter{Event management}

%
% overview of exceptions, interrupts and hardware interrupts
%

\section{Overview of exceptions, interrupts and hardware interrupts}


%
% ---------- header -----------------------------------------------------------
%
% project       kaneton
%
% license       kaneton
%
% file          /home/buckman/kaneton/view/book/architecture/ia32/task.tex
%
% created       matthieu bucchianeri   [sat sep  2 11:41:15 2006]
% updated       matthieu bucchianeri   [sat dec 22 23:11:16 2007]
%

%
% task management
%

\chapter{Task management}

This chapter describes all the IA-32 specific mecanisms used in the
thread, task and scheduler managers.

These mecanisms mainly includes the manipulation of contexts. The
context is the whole state of the microprocessor. The machine
dependent operations consists in creating, maintaining and changing
the contexts.

\newpage

%
% reserving a task
%

\section{Reserving a task} % XXX

%
% preparing a new thread
%

\section{Preparing a new thread} % XXX

%
% scheduling and context switch
%

\section{Thread switching}

The thread's scheduling involves many mecanisms and operations on
contexts. The following sections describes the manipulations of
contexts and the context switching as implemented in kaneton for
IA-32.

%
% context saving and restoring
%

\subsection{Context saving and restoring}

Saving the context is done only on interrupt or exception. Restoring
the context is done only when resuming from such event.

The code of the low-level interrupt and exception handlers must be
generated using one of these four macro functions:

\begin{itemize}
\item \textit{IA32\_EXCEPTION\_PREHANDLER\_CODE}, for exceptions with error
code appended by the microprocessor ;
\item \textit{IA32\_EXCEPTION\_PREHANDLER\_NOCODE}, for exceptions without
any error code ;
\item \textit{IA32\_INTERRUPT\_PREHANDLER}, for interrupts ;
\item \textit{IA32\_IRQ\_PREHANDLER}, for interrupts with PIC aknowlegment ;
\item \textit{IA32\_SYSCALL\_PREHANDLER}, for system calls.
\end{itemize}

The macro functions generates the context management code, which
skeleton is:

\begin{itemize}
\item Switching to the thread's kernel stack (hardware performed) ;
\item Saving the context, \textit{IA32\_SAVE\_CONTEXT()} ;
\item Switching to the kernel interrupt stack ;
\item Calling the higher-level handler ;
\item Switching to the resuming thread's kernel stack ;
\item Restoring the context, \textit{IA32\_RESTORE\_CONTEXT()} ;
\item Returning from interrupt or exception.
\end{itemize}

The \textit{IA32\_SAVE\_CONTEXT} step makes a \textit{pusha} to save
all the general purpose registers and then pushes the data segment
register. Then, it changes the value of the page-directory base
register with kernel page-directory base (stored by the kernel in the
\textit{ia32\_interrupt\_pdbr} variable). Next it sets the kernel
segments (variable \textit{ia32\_interrupt\_ds}). To finish, it saves
the current stack pointer to \textit{ia32\_local\_jump\_stack} before
jumping to the interrupt dedicated stack
(\textit{ia32\_local\_interrupt\_stack}).

\textit{RESTORE\_CONTEXT} does the opposite: it switches back to the
stack \textit{ia32\_local\_jump\_stack}, it pops every registers
pushed on the stack and then switches address space to the one stored
in \textit{ia32\_local\_jump\_pdbr}.

The code of the generated handlers is then placed into the
\textit{.handler} ELF section, which will be mapped in every address
space.

We'll talk about the FPU and SSE contexts in the section
\textbf{Context switching}.

%
% accessing a thread's context
%

\subsection{Accessing a thread's context}

As the execution context of one thread is stored on its kernel stack,
it is impossible to access it directly since this kernel stack is only
mapped into the thread's address space.

Accessing a thread's context requires calling the \textit{as\_read}
and \textit{as\_write} operations.

The function \textit{ia32\_get\_context} retrieve the address of a
thread's kernel stack and uses \textit{as\_read} to read the stored
context. The function \textit{ia32\_set\_context} does the opposite,
except you can specify a mask of registers to update in the thread's
context. This requires to read the context and then to rewrite only
the masked registers, which is slower but easier for the programer.

%
% context switching
%

\subsection{Context switching}

The context of a thread is stored onto the kernel stack of the thread
when an exception or an interrupt occurs. When resuming, the context
at the top of the stack is used to resume execution.

Therefore, the implementation of \textit{sched\_switch} dependent code
consists only in few steps:

\begin{enumerate}
\item Setting the \textit{ia32\_local\_jump\_pdbr} variable to the
PDBR of the resuming thread ;
\item Setting the \textit{ia32\_local\_jump\_stack} variable to the
address of the kernel stack of the resuming thread ;
\item Setting up the TSS to switch stack to the thread's kernel stack
on interrupt ;
\item Copying the permission bitmap of the task into the active TSS
(see \textbf{I/O management} for details about the permission bitmap).
\end{enumerate}

The context switching of FPU and extended registers (SSE and SSE2) is
implemented using IA-32 facilities: when the TS bit of CR0 is set, an
exception is thrown each time a floating or extension instruction is
executed.

Thus, switching from one thread to another sets the TS bit to 1. If
the thread executes no special instruction, the FPU \& extended context
stay unchanged. If the thread executes such instruction, an exception
is thrown. The exception handler (\textit{ia32\_extended\_context\_switch}):

\begin{enumerate}
\item Save the FPU \& extended register set into the \textit{mmx\_context}
thread ;
\item Restore the FPU \& extended context of the currently executing thread ;
\item Clear the TS flag ;
\item Set \textit{mmx\_context} to the current thread id.
\end{enumerate}

\textit{mmx\_context} is a scheduler variable that contains the
identifier of the last thread that executed a FPU or extension
instruction.

Saving and restoring the FPU \& extended context is done using
dedicated instructions: \textit{fsave} and \textit{frstor} when no SSE
extensions are present, \textit{fxsave} and \textit{fxrstor} when SSE
is supported.

%
% address space switching
%

\subsection{Address space switching}

No additional operation is required, since loading the CR3 register
involves flushing the translation caches.

%
% stack switching switching
%

\subsection{Stack switching}

As explained above, the code executed by the kernel on
interrupt/exception runs on a global kernel stack.

XXX creation + TSS

%
% ---------- header -----------------------------------------------------------
%
% project       kaneton
%
% license       kaneton
%
% file          /home/mycure/kaneton/view/lecture/kernels/interrupts/io.tex
%
% created       julien quintard   [wed may 16 19:17:35 2007]
% updated       julien quintard   [wed may 16 19:17:36 2007]
%

\pgfdeclareimage[interpolate=true,width=180pt]
                {io_mapped}
                {figures/io_mapped}

\pgfdeclareimage[interpolate=true,width=200pt]
                {io_port}
                {figures/io_port}

\pgfdeclareimage[interpolate=true,width=320pt]
                {io_hybrid}
                {figures/io_hybrid}

%
%
%

\begin{frame}
  \frametitle{I/O device}

  An I/O device is generally composed of a control register and a data buffer.

  \-

  Example of I/O devices: Hard disk, keyboard, Interrupt Controller, serial
  port\ldots

  \-

  To address the I/O devices, the microprocessor can use one of the 3
  architectures:

  \begin{itemize}
    \item memory-mapped I/O\\
      The hardware device's memory is mapped in RAM.
    \item I/O ports\\
    \item hybrid\\
  \end{itemize}

\end{frame}

%
%
%

\begin{frame}
  \frametitle{Memory-mapped I/O}

  \begin{center}
    \pgfuseimage{io_mapped}
  \end{center}

  {\bf Advantages}\\
  \begin{itemize}
    \item restricted instruction set
    \item No special protection mechanism is needed, since protection is
      already implemented at the memory level.
  \end{itemize}

  {\bf Disadvantages}\\
  \begin{itemize}
    \item Caching. Solution: disable caching (for a page)
  \end{itemize}

\end{frame}

%
%
%

\begin{frame}
  \frametitle{I/O ports}

  \begin{center}
    \pgfuseimage{io_port}
  \end{center}

  {\bf Advantages}\\
  \begin{itemize}
    \item Can accept a large number of devices without affecting main memory.
  \end{itemize}

  {\bf Disadvantages}\\
  \begin{itemize}
    \item Enlarge the instruction set with special read and write instructions.
  \end{itemize}

\end{frame}

%
%
%

\begin{frame}
  \frametitle{Hybrid I/O}

  \begin{center}
    \pgfuseimage{io_hybrid}
  \end{center}

\end{frame}

%
%
%

\begin{frame}
  \frametitle{I/O performing}

  I/O are asynchronous:

  \begin{itemize}
    \item in emission: a hardware device may send data at any time
    \item in reception: the kernel cannot find out when the device is ready to
      receive data
  \end{itemize}

  \-

  There are 2 methods to handle asynchronous events:

  \begin{itemize}
    \item Programmed I/O (polling)
    \item Interrupt-based I/O
  \end{itemize}

\end{frame}

%
%
%

\begin{frame}
  \frametitle{Programmed I/O}



\end{frame}


%
%
%

\begin{frame}
  \frametitle{Interrupt-based I/O}



\end{frame}

%
%
%

\begin{frame}
  \frametitle{Direct Memory Access (DMA)}

  To support DMA, the hardware must have a DMA controller. This controller
  provides the following registers:

  \begin{itemize}
    \item The address register determiness the transfert's memory address.
    \item count register specifies the transfer size.
    \item control register contains the I/O port and the transfert direction.
  \end{itemize}

  \-

  {\bf Advantages}\\
  \begin{itemize}
    \item DMA minimize the microprocessor load
    \item decrease the interrupt rate
  \end{itemize}

  {\bf Disadvantages}\\
  \begin{itemize}
    \item The DMA controller is much slower than the CPU. Also, when the CPU,
      it could assume the transfert faster.
    \item 
  \end{itemize}

\end{frame}

%
% ---------- header -----------------------------------------------------------
%
% project       kaneton
%
% license       kaneton
%
% file          /home/mycure/kaneton/view/book/kaneton/references.tex
%
% created       julien quintard   [tue jun 19 15:36:51 2007]
% updated       julien quintard   [sun dec 16 17:33:44 2007]
%

%
% ---------- references -------------------------------------------------------
%

\chapter{References}
\label{chapter:references}

This chapter contains a listing of the kaneton documents that readers might
be interested in as well.

\newpage

%
% ---------- text -------------------------------------------------------------
%

This document aimed at presenting the kaneton microkernel and its organisation.

The books listed below go further by describing more precisely a specific
component of the kaneton microkernel:

\begin{enumerate}
  \item
    \textbf{The kaneton microkernel :: core}

    \-

    That document focuses on the design and implementation of the kaneton
    microkernel core through the different managers;
  \item
    \textbf{The kaneton microkernel :: \textit{architecture}}

    \-

    Theses documents tackle the implementation of kaneton on the given
    \textit{architecture}. Additionally, these documents contain information
    on platforms coupled with this architecture.

    \-

    Thes documents can be found in two forms, either public or private. Indeed,
    these documents contain information students must not have access to.

    \-

    The architectures currently documented are listed below:

    \begin{itemize}
      \item
	\textbf{ia32}

	\-

	\textit{Platforms}: \textbf{ibm-pc}
    \end{itemize}
\end{enumerate}

%%
%% licence       kaneton licence
%%
%% project       kaneton
%%
%% file          /home/mycure/kaneton/view/papers/kaneton/bibliography.tex
%%
%% created       julien quintard   [mon may  8 18:35:35 2006]
%% updated       julien quintard   [mon may  8 20:38:56 2006]
%%

%
% bibliograpy
%

\chapter{Bibliography}

This chapter contains the bibliography.

%
% text
%

\begin{thebibliography}{0}
  \bibitem{AST-SCO}
    \textbf{Structured Computer Organization};
    by
    \textit{Andrew S. Tanenbaum}
  \bibitem{AST-CN}
    \textbf{Computer Networks};
    by
    \textit{Andrew S. Tanenbaum}
  \bibitem{AST-OSDI}
    \textbf{Operating Systems: Design and Implementation};
    by
    \textit{Andrew S. Tanenbaum, Albert S Woodhull}
  \bibitem{AST-MOS}
    \textbf{Modern Operating Systems};
    by
    \textit{Andrew S. Tanenbaum}
  \bibitem{AST-DOS}
    \textbf{Distributed Operating Systems};
    by
    \textit{Andrew S. Tanenbaum}
  \bibitem{AST-DSPP}
    \textbf{Distributed Systems: Principles and Paradigms};
    by
    \textit{Andrew S. Tanenbaum, Maarten van Steen}
  \bibitem{NAL-DA}
    \textbf{Distributed Algorithms};
    by
    \textit{Nancy A. Lynch}
\end{thebibliography}


\end{document}
